\Needspace{17\baselineskip}
\Article{21}

Zweitens sei $n$ das Quadrat oder eine höhere Potenz einer ungeraden Primzahl $p$, und es werde $n=p^{2\varkappa}q$ gesetzt, wobei $q$ entweder $=1$ oder $=p$ ist.\quad Zunächst ist hier zu beachten, dass für eine beliebige ganze Zahl $\lambda$, die nicht durch $p^\varkappa$ teilbar ist, gilt:

\SourceLeaf{43}

\[
\begin{aligned}
&r^{\lambda\lambda}+r^{(\lambda+p^\varkappa q)^2}+r^{(\lambda+2p^\varkappa q)^2}+r^{(\lambda+3p^\varkappa q)^2}\\&\qquad+\text{usw.}+r^{(\lambda+n-p^\varkappa q)^{3}}\\
&\quad=r^{\lambda\lambda}\{1+r^{2\lambda p^\varkappa q}+r^{4\lambda p^\varkappa q}\\&\qquad+r^{6\lambda p^\varkappa q}+\text{usw.}+r^{2\lambda(n-p^\varkappa q)}\}\\&\quad=\frac{r^{\lambda\lambda}(1-r^{2\lambda n})}{1-r^{2\lambda p^\varkappa q}}=0
\end{aligned}
\]

Man erkennt dann leicht die Gleichung

\[
\begin{aligned}
W=1+r^{p^{2\varkappa}}+r^{4p^{2\varkappa}}+r^{9p^{2\varkappa}}+\text{usw.}+r^{(n-p^\varkappa)^2}
\end{aligned}
\]

Denn die übrigen Glieder der Reihe

\[
\begin{aligned}
1+r+r^4+r^9+\text{usw.}+r^{(n-1)^2}
\end{aligned}
\]

lassen sich in $(p^\varkappa-1)q$ Teilreihen mit jeweils $p^\varkappa$ Gliedern aufteilen, deren Summen durch die soeben angegebene Umformung verschwinden.

Daraus folgt im Fall $q=1$, also wenn $n$ eine Primzahlpotenz mit geradem Exponenten ist:

\[
\begin{aligned}
W=p^\varkappa=+\surd n,\;\text{und somit}\;T=+\surd n,\quad U=0
\end{aligned}
\]

Im anderen Fall dagegen, in dem $q=p$, also $n$ eine Primzahlpotenz mit ungeradem Exponenten ist, setzen wir $r^{p^{2\varkappa}}=\rho$. Dann ist $\rho$ eine eigentliche Wurzel der Gleichung $x^p-1=0$, und zwar $\rho=\cos\frac{k}{p}360^0+i\sin\frac{k}{p}360^0$, und es ergibt sich

\[
\begin{aligned}
W&=1+\rho+\rho^4+\rho^9+\text{usw.}+\rho^{(p^{\varkappa+1}-1)^2}\\&=p^\varkappa(1+\rho+\rho^4+\rho^9+\text{usw.}+\rho^{(p-1)^2})
\end{aligned}
\]

Die Summe der Reihe $1+\rho+\rho^4+\rho^9+\text{usw.}+\rho^{(p-1)^2}$ ist aber durch den vorigen Artikel bestimmt. Daraus folgt unmittelbar

\[
\begin{aligned}
W&=\pm\surd n=T,\quad\text{falls }p\text{ die Form }4\mu+1\text{ hat}\\
W&=\pm i\surd n=iU,\quad\text{falls }p\text{ die Form }4\mu+3\text{ hat}
\end{aligned}
\]

mit dem positiven oder negativen Vorzeichen, je nachdem, ob $k$ ein Rest oder ein Nichtrest von $p$ ist.

\Needspace{12\baselineskip}
\Article{22}

Aus den Ausführungen in den Art. 20 und 21 lässt sich auch leicht der folgende Satz ableiten, der uns weiter unten von erheblichem Nutzen sein wird.\quad Setzen wir

\[
\begin{aligned}
W'=1+r^h+r^{4h}+r^{9h}+\text{usw.}+r^{h(n-1)^2}
\end{aligned}
\]

\SourceLeaf{44}

wobei $h$ eine beliebige, nicht durch $p$ teilbare ganze Zahl bezeichnet, so gilt in dem Fall $n=p$ oder wenn $n$ eine Potenz von $p$ mit ungeradem Exponenten ist:

\[
\begin{aligned}
W'&=W,\;\text{falls }h\text{ ein quadratischer Rest von }p\text{ ist}\\
W'&=-W,\;\text{falls }h\text{ ein quadratischer Nichtrest von }p\text{ ist}
\end{aligned}
\]

Denn offenbar geht $W'$ aus $W$ hervor, wenn für $k$ der Ausdruck $kh$ eingesetzt wird. Im ersten Fall sind aber $k$ und $kh$ hinsichtlich ihrer Eigenschaft als Reste oder Nichtreste von $p$ gleichartig, im zweiten Fall verschiedenartig.

In dem Fall dagegen, dass $n$ eine Potenz von $p$ mit geradem Exponenten ist, gilt offenbar $W'=+\surd n$ und somit stets $W'=W$.

\Needspace{13\baselineskip}
\Article{23}

In den Art. 20, 21 und 22 haben wir ungerade Primzahlen und deren Potenzen betrachtet. Es bleibt daher der Fall, in dem $n$ eine Zweierpotenz ist.

Für $n=2$ ist offenbar $W=1+r=0$.

Für $n=4$ ergibt sich $W=1+r+r^4+r^9=2+2r$. Somit ist $W=2+2i$, wenn $k$ die Form $4\mu+1$ hat, und $W=2-2i$, wenn $k$ die Form $4\mu+3$ hat.

Für $n=8$ haben wir $W=1+r+r^4+r^9+r^{16}+r^{25}+r^{36}+r^{49}=2+4r+2r^4$ $=4r$.\quad Daher ist

\[
\begin{aligned}
W&=(1+i)\surd8,\quad\text{wenn }k\text{ die Form }8\mu+1\text{ hat}\\
W&=(-1+i)\surd8,\quad\text{wenn }k\text{ die Form }8\mu+3\text{ hat}\\
W&=(-1-i)\surd8,\quad\text{wenn }k\text{ die Form }8\mu+5\text{ hat}\\
W&=(1-i)\surd8,\quad\text{wenn }k\text{ die Form }8\mu+7\text{ hat}
\end{aligned}
\]

Ist $n$ eine höhere Zweierpotenz, so setzen wir $n=2^{2\varkappa}q$, wobei $q$ entweder $=1$ oder $=2$ und $\varkappa$ größer als $1$ ist.\quad Zunächst ist hier zu beachten, dass für eine beliebige ganze Zahl $\lambda$, die nicht durch $2^{\varkappa-1}$ teilbar ist, gilt:

\[
\begin{aligned}
&r^{\lambda\lambda}+r^{(\lambda+2^\varkappa q)^2}+r^{(\lambda+2\pdot2^\varkappa q)^2}+r^{(\lambda+3\pdot2^\varkappa q)^2}\\&\qquad+\text{usw.}+r^{(\lambda+n-2^\varkappa q)^2}\\
&\quad=r^{\lambda\lambda}\{1+r^{2^{\varkappa+1}\lambda q}+r^{2\pdot2^{\varkappa+1}\lambda q}\\&\qquad+r^{3\pdot2^{\varkappa+1}\lambda q}+\text{usw.}+r^{(2n-2^{\varkappa+1}q)\lambda}\}\\&\quad=\frac{r^{\lambda\lambda}(1-r^{2\lambda n})}{1-r^{2^{\varkappa+1}\lambda q}}=0
\end{aligned}
\]

Man erkennt dann leicht die Gleichung

\[
\begin{aligned}
W=1+r^{2^{2\varkappa-2}}+r^{4\pdot2^{2\varkappa-2}}+r^{9\pdot2^{2\varkappa-2}}+\text{usw.}+r^{(n-2^{\varkappa-1})^2}
\end{aligned}
\]

\SourceLeaf{45}

Setzen wir $r^{2^{2\varkappa-2}}=\rho$, so ist $\rho$ eine Wurzel der Gleichung $x^{4q}-1=0$, und zwar $\rho=\cos\frac{k}{4q}360^0+i\sin\frac{k}{4q}360^0$; dann ergibt sich

\[
\begin{aligned}
W&=1+\rho+\rho^4+\rho^9+\text{usw.}+\rho^{(2^{\varkappa+1}q-1)^2}\\
&=2^{\varkappa-1}(1+\rho+\rho^4+\rho^9+\text{usw.}+\rho^{(4q-1)^2})
\end{aligned}
\]

\Needspace{14\baselineskip}
Die Summe der Reihe $1+\rho+\rho^4+\rho^9+\text{usw.}+\rho^{(4q-1)^2}$ ist aber durch unsere Ausführungen zu den Fällen $n=4$, $n=8$ bestimmt. Daraus schließen wir,

dass im Fall $q=1$, also wenn $n$ eine Potenz der Zahl $4$ ist,

\[
\begin{aligned}
W&=(1+i)2^\varkappa=(1+i)\surd n,\quad\text{falls }k\text{ die Form }4\mu+1\text{ hat}\\
W&=(1-i)2^\varkappa=(1-i)\surd n,\quad\text{falls }k\text{ die Form }4\mu+3\text{ hat}
\end{aligned}
\]

gilt. Dies sind genau die für $n=4$ angegebenen Formeln;

im anderen Fall dagegen, in dem $q=2$, also $n$ eine Zweierpotenz mit einem ungeraden Exponenten größer als $3$ ist, gilt

\[
\begin{aligned}
W&=(1+i)2^\varkappa\surd2=(1+i)\surd n,\quad\text{falls }k\text{ die Form }8\mu+1\text{ hat}\\
W&=(-1+i)2^\varkappa\surd2=(-1+i)\surd n,\quad\text{falls }k\text{ die Form }8\mu+3\text{ hat}\\
W&=(-1-i)2^\varkappa\surd2=(-1-i)\surd n,\quad\text{falls }k\text{ die Form }8\mu+5\text{ hat}\\
W&=(1-i)2^\varkappa\surd2=(1-i)\surd n,\quad\text{falls }k\text{ die Form }8\mu+7\text{ hat}
\end{aligned}
\]

was ebenfalls vollständig mit den für $n=8$ angegebenen Formeln übereinstimmt.

\Needspace{12\baselineskip}
\Article{24}

Auch hier wird es sich lohnen, das Verhältnis der Summe der Reihe

\[
\begin{aligned}
W'=1+r^h+r^{4h}+r^{9h}+\text{usw.}+r^{h(n-1)^2}
\end{aligned}
\]

zu $W$ zu bestimmen, wobei $h$ eine beliebige ungerade ganze Zahl bezeichnet.\quad Da $W'$ aus $W$ durch Ersetzen von $k$ durch $kh$ hervorgeht, hängt der Wert von $W'$ ebenso von der Form der Zahl $kh$ ab, wie $W$ von der Form von $k$ abhängt.\quad Setzen wir $\frac{W'}{W}=l$, so gilt offenbar

I.\quad im Fall $n=4$ oder einer höheren Zweierpotenz mit geradem Exponenten

\[
\begin{aligned}
l&=1,\quad\text{falls }h\text{ die Form }4\mu+1\text{ hat}\\
l&=-i,\quad\text{falls }h\text{ die Form }4\mu+3\text{ hat},\;\text{und }k\text{ die Form }4\mu+1\text{ hat}\\
l&=+i,\quad\text{falls }h\text{ die Form }4\mu+3\text{ hat},\;\text{und }k\text{ dieselbe Form hat.}
\end{aligned}
\]

\SourceLeaf{46}

\Needspace{15\baselineskip}
II.\quad im Fall $n=8$ oder einer höheren Zweierpotenz mit ungeradem Exponenten

\[
\begin{aligned}
l&=1,\ \text{falls}\  h\text{ die Form }8\mu+1\text{ hat},\\
l&=-1,\ \text{falls}\  h\text{ die Form }8\mu+5\text{ hat},\\
l&=+i,\ \text{falls entweder }h\text{ die Form }8\mu+3\text{ hat},\;\text{und }k\text{ die Form }4\mu+1\text{ hat},\\
&\phantom{=+i,}\text{oder }h\text{ die Form }8\mu+7\text{ hat},\;\text{und }k\text{ die Form }4\mu+3\text{ hat},\\
l&=-i,\ \text{falls entweder }h\text{ die Form }8\mu+3\text{ hat},\;\text{und }k\text{ die Form }4\mu+3\text{ hat},\\
&\phantom{=+i,}\text{oder }h\text{ die Form }8\mu+7\text{ hat},\;\text{und }k\text{ die Form }4\mu+1\text{ hat}.
\end{aligned}
\]

Durch die vorigen Artikel ist die Bestimmung der Summe $W$ für die Fälle, in denen $n$ eine Primzahl oder eine Primzahlpotenz ist, vollständig abgeschlossen. Es bleibt daher, auch die Fälle zu erledigen, in denen $n$ aus mehreren Primzahlen zusammengesetzt ist; hierzu wird uns der folgende Satz den Weg bahnen:
